% Focused draft: general method; one conveyor-based virtual testbed.
\section{Introduction}
% [claim:C102] [evidence:E101,E108,E109,E112,E113]
Cyber-physical production systems execute operations whose progress is often represented by discrete events: a workpiece crosses a sensing point, an actuator reaches an endpoint, or an axis reports arrival. A monitoring decision must relate the received event to the operation that was expected and to the time at which the observation was generated. Digital-twin research distinguishes model representation from actual physical--virtual integration~\cite{Kritzinger2018}; production-system studies emphasize sensing, acquisition, and the services built on that information~\cite{Negri2017,Uhlemann2017,Cimino2019}. Jones et al.~\cite{Jones2020} further identify states, connections, and twinning rate among characteristics of a twin. These foundations motivate treating an observation's availability as part of a monitoring method rather than assuming that a simulated state is a measurement.

% [claim:C103] [evidence:E102,E103,E110,E111,E114,E115]
Communication delay and loss are established concerns in networked control~\cite{Hespanha2007}, and age of information provides a time-based description of freshness~\cite{Yates2021}. Reviews discuss physical and learned models for digital twins~\cite{Rasheed2020,Fuller2020}. Neither a more detailed model nor a lower prediction error alone establishes that its output is suitable for a particular decision. Shao et al.~\cite{Shao2023} place verification, validation, and uncertainty assessment when assessing the credibility of a manufacturing twin. Lu et al.~\cite{Lu2024} examine how quality issues propagate through the development of digital twins. Their work supports evaluating the data and model basis of a service; it does not validate the thresholds used in the present testbed.

% [claim:C104] [evidence:E104,E105,E009,E012]
Prior applications use twins to generate training data for manufacturing fault diagnosis~\cite{Ademujimi2022} and to assess how manufacturing lines can be reconfigured under disturbances~\cite{Fu2025}. The present paper asks a narrower question: \emph{when can a predicted event window support an interpretable finding of inconsistent process progress under incomplete observations, and does a learned residual improve that finding over a simple correction?} Prediction error, warning timeliness, evidence availability, and fault attribution are evaluated as distinct outcomes. A missing transition is not automatically assigned to a uniquely identified physical fault.

% [claim:C105] [evidence:E009,E012,E013]
We propose a process-independent event-window monitoring interface with explicit eligibility conditions, and implement a reusable binary-event kernel behind a process-specific prediction adapter. A matched virtual comparison evaluates a pulse-based predictor, a training-mean correction, and Ridge residual learning. The cell is a testbed for this formulation, not the definition of the method. The intended abstraction includes other predictable binary production events, but the empirical evaluation concerns one conveyor-based system only. This is an exploratory methodological case study, not a systematic literature review, a claim of general algorithmic novelty, or a demonstration of autonomous production recovery.

\section{Problem Formulation and Applicability}
% [claim:C106] [evidence:E013]
Consider a production operation with a past, observable anchor at time $t_a$ and a subsequent binary completion event $e$. A predictor $P_m$ uses admissible observations and issued commands up to $t_a$ to produce an event-time estimate $\widehat t_m$ and a window $I_m=[\ell_m,u_m]$. The anchor need not be a conveyor pulse: it can be an observed operation start or another validated transition. The predictor and its calibration are specific to that event and operating domain.

An observation is represented as
\begin{equation}
 o_k=(s_k,r_k,n_k,y_k,q_k),\qquad A_k=t_k-s_k,
\end{equation}
where $s_k$ is source time, $r_k$ reception time, $n_k$ sequence number, $y_k$ the binary event value, $q_k$ quality metadata, and $t_k$ decision time. Define eligibility as the conjunction
\begin{equation}
 G_k=G_{\mathrm{clock}}\land G_{\mathrm{age}}\land G_{\mathrm{quality}}
 \land G_{\mathrm{history}}\land G_{\mathrm{domain}}.
\end{equation}
These terms require compatible clocks with $s_k\leq r_k\leq t_k$, age within the configured limit, resolved signal quality, sufficient contiguous edge history, and applicability of the prediction and calibration. Missing metadata cannot silently be replaced by an assumed healthy value. Here, clock compatibility is provided by the shared simulation clock; physical clock alignment is not implemented.

The formulation requires a predictable event, a causal anchor, a calibrated operating domain, and an explicit relation between local commands and the observed operation. It is unsuitable when the expected event cannot be predicted from available information, when unknown command changes invalidate the forecast, or when competing operations cannot be distinguished. One monitor instance covers one operation; overlapping products and automatic reassignment of event identity are outside the evaluated implementation.

\section{Monitoring Method}
\subsection{Separation of Prediction and Event Evaluation}
% [claim:C107] [evidence:E009,E012,E013]
The predictor supplies a window; the event kernel consumes that window and eligible binary evidence. The kernel contains no conveyor coordinates, actuator equations, robot geometry, PLC addresses, or MuJoCo calls. Such information belongs in the process adapter. Consequently, changing equipment does not require changing the event comparison logic, but it does require establishing a new observation mapping, predictor, calibration, and eligibility assumptions. Interface reuse is not evidence that an existing calibrated model transfers between machines.

For a point predictor and an empirical residual radius, the implemented window is
\begin{equation}
 I_m=[\max(t_a,\widehat t_m-r_m),\widehat t_m+r_m],\quad
 r_m=\max_{i\in\mathcal D_{\mathrm{cal}}}|t_{e,i}-\widehat t_{m,i}|.
\end{equation}
The radius uses a separate normal calibration set. It is an observed maximum on that set, not a confidence bound or a calibrated probability of a future event. This choice is deliberately simple so that differences in predictor error and window width remain inspectable.

\subsection{Findings, Abstention, and Evidence Attribution}
% [claim:C108] [evidence:E012,E013]
For a resolved rising edge at source time $s_k$, a finding is early if $s_k<\ell_m$, late if $s_k>u_m$, and otherwise consistent with the window. If no target edge has been seen, missing progress is flagged only after the configured number of consecutive eligible samples satisfy $s_k>u_m$. The case study uses two samples. A high value at the first eligible observation cannot retrospectively establish the time of its rising edge and is labeled unresolved.

When $G_k$ is false, the monitor reports insufficient evidence rather than a healthy process or an inferred missing event. Lost edge history invalidates the current kernel instance; later reception alone does not restore its missing history. The adapter can select a separate engineering-bound predictor when an empirical model is outside its domain, but that result retains a different attribution. In particular, an engineering fallback alert must not be credited to a learned predictor.

The source-time condition matters under delay: lack of a received edge at decision time does not establish absence through that time. An old packet only reports the process through its source timestamp. The monitored outcome is therefore an inconsistency of expected progress with admissible evidence, not necessarily a classified sensor, communication, or drive fault. No monitor result has actuator or PLC output authority in this study.

\subsection{Causal Learning and a Matched Simple Reference}
% [claim:C109] [evidence:E009,E012]
A learned residual can supplement a nominal event-time estimate:
\begin{equation}
 \widehat t_{\mathrm{learned}}=\widehat t_{\mathrm{base}}+b+w^T\widetilde\phi.
\end{equation}
Features $\phi$ must be available before the target event; scaling, fitting, model selection, and envelope calibration use their assigned data partitions. The learned alternative is compared with the same baseline plus a training-mean residual. This reference tests whether apparent learning benefit is principally a systematic offset correction. A learned model is not selected merely because its point-prediction error is lower: the resulting window, false findings, rejection rate, and warning time also matter.

\section{Virtual Testbed and Experimental Protocol}
\subsection{Process-Specific Adapter}
% [claim:C110] [evidence:E003,E009,E012,E106]
A conveyor-based production cell supplies a single virtual case study. MuJoCo provides the multibody/contact simulation environment~\cite{Todorov2012}; the monitoring inputs are separated from simulator pose and hidden fault labels. The broader cell includes a pusher and gantry, but the reported experiment isolates one workpiece moving between two optical sensing locations. It does not evaluate the complete pick-and-place cycle. CAD V5 and the functional MuJoCo baseline remain different geometries and are not presented as synchronized representations of a calibrated physical machine.

The adapter uses the completed upstream occupancy pulse as its anchor. If $L$ is nominal workpiece length, $D_p$ commanded travel integrated during the pulse, $v$ constant issued speed, and $d$ the remaining event-to-event travel, then
\begin{equation}
 \widehat g=L/D_p,\qquad \widehat t_{\mathrm{pulse}}=t_a+d/(v\widehat g).
\end{equation}
The Ridge features are pulse duration, issued speed, and their product, frozen at pulse completion. Labels are sampled downstream optical-edge timestamps, not simulator positions. Only constant forward command segments are admitted to this empirical comparison. Source sampling is 20~ms; preliminary workpiece length is 60~mm. Neither number constitutes a new physical calibration campaign.

The adapter checks the training feature box and the stored calibration context. Changes in workpiece length, sensor spacing, sampling period, or command history prevent reuse of the old empirical window. An engineering pulse-bound fallback remains separate; invalid quality or lost history suppresses model findings. These checks do not establish complete domain validity or a physical risk guarantee.

\subsection{Data Partitions and Frozen Comparisons}
% [claim:C111] [evidence:E009]
The deterministic normal campaign contains 61 segments: 32 training, eight validation, eight error-envelope calibration, nine normal test, and four shifted-condition segments. Speed--gain--friction tuples do not overlap between partitions. However, the admitted observation--command pairs are distinct only 16, four, four, nine, and four times, respectively, because the friction variants in the first three partitions produce identical recordings. No pair overlaps between partitions. These designed conditions and duplicates are not independent factory observations.

Ridge penalties $0.01$, $0.1$, $1$, and $10$ are selected by validation mean absolute error (MAE); scaling and weights remain fitted on training rows. The model is sealed before calibration and test generation. A constant-bias reference was added after the initial pilot and is disclosed as post hoc; its offset uses training rows only. Ridge features, penalty selection, and weights were not revised from test outcomes. A first generation attempt failed before a usable training segment because the fixture inherited an incompatible degraded-speed setting; the local fixture correction and failed attempt remain documented.

% [claim:C112] [evidence:E012,E013]
Monitoring evaluation retrospectively reuses the nine normal test and four shifted recordings and adds missing-target feedback and a stalled drive for each normal test condition: 31 recordings in total. Three methods and five scheduled delivery cases produce 465 replays. Cases are immediate reception, 40~ms delay, 120~ms stale delivery, a brief post-pulse gap, and unavailable target-signal quality. Model weights and separately calibrated radii remain fixed. This is exploratory replay, not a new independent confirmation set or a measurement of physical network performance.

Outcomes are point-prediction error, findings in normal segments, findings in faulty segments, evidence rejection, fallback attribution, and paired warning time. A recording ending before its required confirmation window is censored rather than silently classified as healthy or a proven missed detection. Fault injection time and simulator state are evaluator-only information. Replays of the same recording are matched software conditions, not extra independent trials.

\section{Results}
\subsection{Prediction Error Is Not Warning Timeliness}
% [claim:C113] [evidence:E009,E012]
\begin{table}[t]
\caption{Sampled optical-edge arrival-time prediction: held-out normal test ($n=9$) and shifted conditions ($n=4$). MAE in ms; not physical timing accuracy.}
\label{tab:eta-ml}
\centering\small
\begin{tabular}{lrr}
\toprule
Method & Test MAE & Shift MAE\\
\midrule
Nominal command kinematics & 105.7 & 330.0\\
Upstream pulse kinematics & 98.5 & 83.3\\
Pulse plus train-mean bias & 18.1 & 28.5\\
Pulse plus Ridge residual & 15.4 & 38.3\\
\bottomrule
\end{tabular}
\par\smallskip\footnotesize Constant-bias reference added after the initial pilot; its value uses training rows only.
\end{table}

On the nine normal test segments, Ridge MAE is 15.4~ms versus 18.1~ms for constant correction and 98.5~ms for the uncorrected pulse estimate. Its incremental advantage over the simple corrected reference is only about 2.7~ms, without a significance claim. Under the four shifted conditions, Ridge is worse than constant correction (38.3 versus 28.5~ms); all four lie outside the training feature box. Sampling and the small designed grid preclude inferring physical sub-sample timing accuracy.

The calibration radii are 140.0~ms for the pulse estimate, 42.9~ms for constant correction, and 65.0~ms for Ridge. Thus the learned point estimate has lower normal-test MAE but a wider event window than the corrected reference. The eight calibration rows contain only four distinct admitted recording pairs, and these radii are not statistical coverage guarantees.

\subsection{Findings and Unavailable Evidence}
% [claim:C114] [evidence:E012,E013]
\begin{table}[t]
\caption{Advisory ETA findings in fixed replay conditions; fault denominator combines nine missing-target and nine stalled-drive segments, which form nine identical input pairs.}
\label{tab:eta-monitor}
\centering\small
\begin{tabular}{lrrrr}
\toprule
Method & Radius & Normal alerts & Fault alerts & Fault alerts\\
 & (ms) & no delay (/9) & no delay (/18) & 40 ms (/18)\\
\midrule
Pulse & 140.0 & 0 & 18 & 16\\
Constant bias & 42.9 & 0 & 18 & 18\\
Ridge & 65.0 & 0 & 18 & 18\\
\bottomrule
\end{tabular}
\par\smallskip\footnotesize Two delayed pulse replays end before a late-edge confirmation window. Data rejection and engineering fallback are not credited as learned-model detections.
\end{table}

Under immediate reception, each method flags all 18 fault segments and none of the nine normal segments. Under 40~ms delay, both corrected alternatives flag 18 of 18; the uncorrected pulse method flags 16. The two remaining pulse replays end before a late-edge confirmation window and therefore lack sufficient observation for a missed-detection conclusion. Under usable delivery, the four shifted normal recordings use engineering fallback and produce no findings; these are not learned-model successes.

Relative to constant correction, Ridge warns at the same time in two immediate fault replays and 20--40~ms later in 16; it is never earlier. The 40~ms delivery case preserves that ordering. This descriptive result illustrates why improvement in a point-prediction metric cannot be substituted for improvement in a decision metric.

All stale-delivery inputs are rejected before a usable fit. Post-pulse gaps and target-quality loss stop admissible evaluation after initial prediction. Zero model findings in these cases reflects insufficient evidence, not correct recognition of healthy operation. All nine missing-sensor/stalled-drive pairs have identical admitted observation and issued-command hashes. The findings consequently cannot uniquely attribute the physical cause, and the 18 fault segments do not represent 18 independent observation patterns.

\subsection{Fresh Conditions: Clock Choice and Age-Gate Ablation}
% [claim:C117] [evidence:E014]
A subsequent protocol was committed before outcome inspection (source revision \texttt{5a39fe5}). It fixes three new commanded speeds, two gains, and two initial longitudinal positions, giving 12 conditions disjoint in speed--gain pairs from the earlier learning grid. Normal, missing-target, and stalled-drive segments produce 36 recordings. Four clock/age policies and five delivery cases yield 720 matched replays, not independent trials. The predictor is the frozen constant correction; no new training or threshold tuning is performed. The zero-normal-finding budget is an observed-count target, not a bound on a population error probability.

The factors are comparison time (source or receiver) and age rejection (on or off). All alternatives retain quality, edge-history, and domain checks and use identical event windows whenever eligible. The receiver-time alternative intentionally compares a source-anchored window with reception time; it is a clock-use ablation, not a well-designed traditional PLC reference. The age-off variant changes only the age limit in a private receiver configuration. Queued packets are drained after source recording ends without inventing new source observations.
\begin{table}[t]
\caption{Fresh-condition clock ablation: normal finding counts and eligibility among 12 planned normal recordings. All rows use the same frozen constant correction.}
\label{tab:clock}
\centering\small
\begin{tabular}{rlrrr}
\toprule
Delay & Age gate & Source findings & Receiver findings & Eligible\\
(ms) & & (/12) & (/12) & (/12)\\
\midrule
0 & On & 2 & 2 & 6\\
40 & On & 2 & 3 & 6\\
80 & On & 2 & 5 & 6\\
120 & On & 0 & 0 & 0\\
120 & Off & 2 & 5 & 6\\
\bottomrule
\end{tabular}
\par\smallskip\footnotesize Six recordings lack complete upstream pulses. The zero-finding target already fails without delay; zero eligible records under stale rejection is not healthy-process evidence.
\end{table}


Six normal recordings start with the upstream sensor already active, preventing recovery of a complete occupancy pulse. Only six of the 12 normal conditions admit an empirical prediction under immediate reception. The source-time alternative nevertheless produces two findings on those six normal recordings even without delay; the frozen calibration therefore fails the stated zero-finding target on fresh conditions. It is not re-fitted on these outcomes.

At 80~ms delay, source and receiver comparison produce two and five normal findings on the same six eligible records, respectively. Under immediate, 40~ms, and 80~ms delivery, each alternative flags 12 of 24 fault segments, with the remaining 12 lacking an admissible pulse. Those 24 segments form 12 identical missing-sensor/stalled-drive input pairs; the 36 recordings contain 24 distinct observation--command pairs. In the eligible delayed fault cases, receiver-time missing findings can occur before the latest delivered source sample passes the window endpoint. This identifies a clock-interpretation error without establishing unique fault causes.

At 120~ms delay the 100~ms age gate rejects all inputs; age-off source comparison retains six eligible normal records and two findings, whereas age-off receiver comparison produces five. Zero findings under age rejection is abstention, not fulfillment of a useful diagnostic objective. Post-pulse gaps invalidate the histories in all policies. Thus time interpretation avoids additional delay-induced findings in this fixture, but neither it nor a rejection rule repairs incomplete anchors or inadequate calibration.

\subsection{Post-Hoc Sampling and Motion-Model Audit}
% [claim:C118] [evidence:E015]
A diagnostic audit separates an ideal sampling bound from actual model validity. With constant physical speed, known aligned length $L$, ideal point beams, and sampled rising, falling, and target edges delayed by $\epsilon_r,\epsilon_f,\epsilon_e\in[0,T_s]$, the pulse estimate satisfies
\begin{equation}
 \begin{aligned}
 \widetilde t_e-\widehat t_e&=\epsilon_e-(1+d/L)\epsilon_f+(d/L)\epsilon_r,\\
 |\widetilde t_e-\widehat t_e|&\leq(1+d/L)T_s.
 \end{aligned}
\end{equation}
Here $d$ is remaining travel and $T_s$ the sampling period. The case-study values give an 86.7~ms ideal radius. Only two of six complete-pulse normal records lie within it. Evaluator-only position traces show that mean speed after the pulse exceeds mean speed during the pulse in all six, by 5.9--7.9\%. This supports a violation of constant-speed assumptions; it does not isolate every geometric or dynamic error. The audit is post hoc and introduces no new calibrated predictor or repaired false-finding result. The offline feature extractor is also corrected to reject an initially high upstream value as an unobserved pulse start. The original 61 training-study results reproduce unchanged after this boundary fix.

\section{Generality, Limitations, and Focused Development}
% [claim:C115] [evidence:E012,E013,E014,E015]
The general element is the separation of a causal event predictor, evidence eligibility, and event-window evaluation. A cylinder-endpoint adapter could anchor at an observed stroke start; an axis-arrival adapter could use a validated motion start and commanded trajectory. Each would need its own event semantics, feedback mapping, dynamic assumptions, and calibration. Their mention describes interface applicability, not validation on those machines. Software contract checks with renamed event labels do not constitute additional physical or simulated-device experiments.

The packaged runtime selects the monitoring method and its MuJoCo case-study dependencies. Full-cell reconfiguration, CAD authoring, graphics, reinforcement learning, and whole-cell online contact forecasting are outside the active paper scope. Historical code, failed and adverse results, parameterized profiles, and the last user-accepted image remain preserved. Scope reduction does not convert earlier timing or recovery limitations into successful results.

The main evidence limitations are one virtual process, deterministic and partly duplicated conditions, reuse of previously inspected normal test recordings, a post hoc simple reference, assumed rather than physically identified dynamics, and empirical calibration maxima. The study establishes neither cross-machine transfer nor live clock synchronization, diagnostic accuracy, hard real-time execution, production availability, or successful operation after loss of required feedback. Model/source hashes and reproducibility checks support traceability, not scientific novelty or a guarantee of acceptance by a venue.

The new fixed-condition ablation isolates clock choice and age rejection, while exposing failure of the old calibration on fresh normal records. The post-hoc motion audit further motivates checking startup/transient dynamics before selecting new calibration data. Next work should cover these effects, reserve another unseen evaluation set, evaluate a second event adapter, and estimate physical model and measurement uncertainty when measurements become available. Until those steps are completed, the general interface and single-testbed results must remain separate claims.

\section{Physical Validation}
\label{sec:physical}
\noindent\textit{Reserved: no physical measurements or validation results are reported.}
\subsection{Physical Setup and Calibration}
\vspace{8mm}
\subsection{Live Synchronization and Network Measurements}
\vspace{8mm}
\subsection{Physical Event-Monitoring Trials}
\vspace{8mm}

\section{Conclusion}
% [claim:C116] [evidence:E009,E012,E013,E014]
A network-aware predictive monitoring framework has been formulated around binary production events and implemented with a process-independent comparison kernel. One conveyor-based MuJoCo testbed demonstrates the distinction between a model mismatch, unavailable evidence, and ambiguous fault attribution. Learned residual correction improves normal-test event-time MAE slightly over a constant correction but does not improve warning timeliness and performs worse under shifted conditions. A later fixed-condition ablation shows fewer delay-induced normal findings with source-time comparison, but the old calibration still violates its normal-finding target and incomplete pulses limit availability. The supported outcome is an inspectable method and bounded comparisons, not verified cross-machine resilience or autonomous PLC reconfiguration.

\section*{Declarations Awaiting Author Completion}
\noindent Authors and affiliations: \hrulefill\par\medskip
\noindent Funding and competing interests: \hrulefill\par\medskip
\noindent Author-approved AI-use disclosure: \hrulefill
